\documentclass[11pt]{article}
\usepackage[utf8]{inputenc}
\usepackage[T1]{fontenc}
\usepackage[brazil]{babel}
\usepackage{lmodern}
\usepackage{amsmath,amssymb}
\usepackage{hyperref}
\usepackage{mathptmx}

% CONFIGURAÇÃO DE TEOREMAS
\usepackage[thmmarks]{ntheorem}

\theoremstyle{plain}
\theoremheaderfont{\normalfont\bfseries}
\theorembodyfont{\normalfont}
\theoremseparator{:}
\theorempreskip{\topsep}
\theorempostskip{\topsep}

\theoremsymbol{\ensuremath{\bigcirc}}
\newtheorem{definicao}{Definição}[section]

\theoremsymbol{\ensuremath{\triangle}}
\newtheorem{nota}[definicao]{Nota}

\theoremsymbol{}
\newtheorem{proposicao}[definicao]{Proposição}

\theoremsymbol{}
\newtheorem{lema}[definicao]{Lema}

\theoremsymbol{\ensuremath{\square}}
\newtheorem{proplivre}[definicao]{Proposição}

\theoremstyle{nonumberplain}
\theoremheaderfont{\normalfont\bfseries}
\theorembodyfont{\normalfont}
\theoremsymbol{\ensuremath{\blacksquare}}
\newtheorem{demonstracao}[definicao]{Demonstração}

% Macros Gerais
\newcommand{\naturais}{\mathbb{N}}
\newcommand{\Z}{\mathbb{Z}}
\newcommand{\R}{\mathbb{R}}
\newcommand{\C}{\mathbb{C}}
\newcommand{\inv}[1]{#1^{-1}}
\newcommand{\set}[1]{\{#1\}}
\newcommand{\braket}[1]{\langle #1 \rangle}
\newcommand{\sseq}{\subseteq}

\title{Autômatos}
\author{Gian Carlo}
\date{\today}

\begin{document}
\maketitle

\section{Aula 2}

\begin{definicao}
    Um $a \in A$ é dito divisor de zero se existe $b$ tal que $ab = 0_A$.
\end{definicao}

\begin{definicao}
    Seja $S \subseteq A$ um subconjunto, denotamos por $\braket S$ o ideal geral pelos elementos de $S$ por somas e multiplicações finitas: \begin{equation}
        i \in \braket S \implies i = \sum_{i = 1}^{n} a_is_ib_i
    \end{equation}
    para $a_i, b_i \in A$ e $s_i \in S$.
\end{definicao}

\begin{nota}
    Temos que $\braket S$ é um subgrupo abeliano de $(A, +)$ (verificar).
\end{nota}

\begin{proposicao}
    $\braket S$ é o menor ($\subseteq$) ideal em $A$ contendo $S$.
\end{proposicao}

\begin{demonstracao}
    Suponha que $I$ seja ideal de $A$ e que $S \sseq I$. Como $I$ é ideal, $aIb \sseq I$. Como $S \sseq I$, $aSb \sseq aIb \sseq I$. Como todo ideal é subgrupo em $+$ do anel, temos que \begin{equation}
        \sum^n aSb \sseq I,
    \end{equation} ou seja, $\braket S \sseq I$.
\end{demonstracao}

\begin{proposicao}
    Para uma cadeia de ideais \begin{equation}
        ... \subseteq I_{i-1} \subseteq I_i \subseteq I_{i+1} \subseteq ...
    \end{equation}
    temos que $I = \bigcup I_i$ é também um ideal em $A$.
\end{proposicao}

\begin{demonstracao}
    Se $x \in I$, existe $I_p$ tal que $x \in I_p$. Daí, $axb \in I_p$ e finalmente $axb \in I$, para $a, b \in A$ quaisquer.
\end{demonstracao}

\begin{proplivre}[Lema de Zorn]
    Para um conjunto não vazio com uma relação de ordem parcial, se toda cadeia ordenada não vazia tem uma cota superior então existe um elemento maximal $m$.
\end{proplivre}

\begin{definicao}
    Um ideal próprio $M \sseq A$ é dito \textbf{maximal} se não existe outro ideal próprio $M \neq I$ tal que $M \sseq I \sseq A$.
\end{definicao}

\begin{nota}
    Em $\R$, $\C$ e $\Z_p$, o único ideal é $\set 0$ que é portanto o ideal maximal.
\end{nota}

\begin{proposicao}
    Em um anel unitário $A$, todo ideal próprio $I$ é contido num ideal maximal.
\end{proposicao}

\begin{demonstracao}
    Considere uma cadeia de ideais contendo $I$: \begin{equation}
        I \sseq J_1 \sseq J_2 \sseq ... \sseq J_i \sseq ...
    \end{equation}
    Temos que $J = \cup J_i$ é uma cota superior. Considere o conjunto $S$ de ideais que contém $I$. Este conjunto é parcialmente ordenado, e pelo lema de Zorn existe pelo menos um elemento maximal $M$ em $S$.
\end{demonstracao}

\begin{nota}
    Em geral, para $I, J$ ideais de um $A$, $I \cup J$ não é um ideal mas é usado para denotar o ideal $\braket{I \cup J}$.
\end{nota}

\begin{proposicao}
    Seja $\phi : A \to B$ um homomorfismo de anéis.

    \begin{enumerate}
        \item Se $R \sseq A$ é um subanel, então $\phi(R)$ é subanel de $B$;
        \item Se $S \sseq B$ é subanel, então $\inv \phi(S)$ é subanel de $A$;
        \item Se $I \sseq A$ é um ideal e $\phi$ é sobrejetora, $\phi(I)$ é um ideal de $B$; e
        \item Se $J \sseq B$ é um ideal então $\inv \phi (I)$ é um ideal em $A$.
    \end{enumerate}
\end{proposicao}

\begin{demonstracao}
    \begin{enumerate}
        \item Seja $R \sseq A$ e $x, y \in R$. Como $\phi(x + y) = \phi(x) + \phi(y)$, e $\phi(x), \phi(y) \in \phi(R)$, temos que $\phi(x + y) \in \phi(R)$. Ocorre a mesma coisa para a multiplicação. Isto nos diz que $\phi(R)$ é fechada pelas operações de $B$, de onde vemos que é subanel.
        \item Sejam $a, b \in \inv \phi (S)$. Como $\phi(a) \in S$ e $\phi(b) \in S$, e $S$ é subanel, $\phi(a) + \phi(b) \in S$ e assim $\phi(a + b) \in S$, e logo $a + b \in \inv \phi (S)$. Analogamente para o produto.
        \item Seja $I \sseq A$ um ideal e $\phi$ sobrejetora. Tome $b \in B$ e $r \in \phi(I)$. Obtemos $a \in A$ com $b = \phi(a)$ e $s \in I$ com $r = \phi(s)$. Assim, \begin{equation}
            br = \phi(a)\phi(s) = \phi(as)
        \end{equation}
        e como $as \in I$, $\phi(as) \in \phi(I)$ e assim $\phi(I)$ é um ideal.
    \end{enumerate}
\end{demonstracao}

\begin{definicao}[Anel Quociente]
    Seja $A$ um anel e $I$ um ideal. Como $I$ é subgrupo abeliano de $(A, +)$, podemos formar as classes laterais. Denotaremos por \begin{equation}
        [x] = x + I = \set {y \in A : x + (-y) \in I}.
    \end{equation}

    Seja $R/I$ o conjunto das classes laterais. Definindo a soma e a multiplicação como \begin{equation}
        [x] + [y] = [x + y]
    \end{equation}
    e \begin{equation}
        [x][y] = [xy]
    \end{equation}
    obtemos uma estrutura de anel. A tripla $(R/I, +, \cdot)$ é o \textbf{anel quociente}.
\end{definicao}

\begin{nota}
    Por exemplo, considere $\C[[x]]$ como o anel das séries formais (não necessariamente convergentes). Como um elemento qualquer tem a forma $\sum a_i x^i$, o anel quociente \begin{equation}
        \C[[x]]/\braket{x^n}
    \end{equation}
    tem os elementos da forma $a_0 + a_1x + ... + a_{n-1}x^{n-1}$.

    $R[x]/\braket{x^2 + 1}$ é isomorfo a $\C$. Para isso, observamos que $x^2 \notin I$ e $-1 \notin I$, porém $x^2 \equiv -1$ já que $x^2 + (-1) \in I$. Partindo da equação $x^2 \equiv -1$, obtemos que o monômio $x$ fica com a estrutura do número $i \in \C$, e que as classes de equivalência ficam na forma $a + bx \cong a + bi$ para $a, b \in \R$.
\end{nota}

\begin{proposicao}
    Seja $I \sseq A$ um ideal. Todo subanel de $A/I$ é o quociente $S/I$ de um único subanel $S \sseq A$ que contém $I$. Todo ideal de $R/I$ é o quociente $J/I$ de um único ideal $J \sseq A$ que contém $I$.
\end{proposicao}

\begin{demonstracao}
    Considere $p : A \to A/I$. Seja $T \sseq A/I$ subanel. Temos que $S = \inv p (T)$ é um subanel de $A$. Se $S'$ é outro subanel de $A$ tal que $p(S') = T$, então $x \in S' \implies p(x) \in T \implies x \in \inv p (T) = S$, ou seja, $S' \sseq S$. Temos que $I \sseq S$ pois $\inv p (I) = \ker p$.

    Seja $K \in R/I$ um ideal. Como $p$ é sobrejetora, então $J = \inv p (K)$ é um ideal.
\end{demonstracao}

\begin{nota}
    Denotaremos por $\pi : A \to A/I$ o homomorfismo que projeta $A$ nas classes de $A/I$.
\end{nota}

\begin{proposicao}[Teorema da Fatorização]
    Se $\phi : A \to B$ é um homomorfismo de anéis tal que $\ker \phi \sseq I$ para algum ideal $I$ de $A$, então existe um homomorfismo $\psi : A/I \to B$ tal que $\psi \circ \pi = \phi$.
\end{proposicao}

\begin{demonstracao}
    Defina $\psi(x + I) = \phi(x)$. Se $x + I = y + I$, temos que $\phi(x) = \phi(y) \implies \phi(x + -y) = 0 \implies x + -y \in \ker \phi$. Daí, $I \sseq \ker \phi$.

    Que $\psi$ é homomorfismo temos de sua definição por $\phi$.
\end{demonstracao}

\begin{proposicao}[Teorema do Isomorfismo]
    Seja $\phi : A \to B$ um homomorfismo de anéis. Então $A/\ker \phi \cong Im(\phi)$, ou seja, existe $\theta : A/\ker \phi \to Im (\phi)$ único tal que o diagrama (\textit{diagrama à fazer}) comuta.
\end{proposicao}

\begin{demonstracao}
    Obtemos $\theta : A/\ker \phi \to Im(\phi)$ fazendo $\ker \phi = I$ na proposição anterior. Para $b \in Im(\phi)$, temos $a \in A$ tal que $\theta (a + \ker \phi) = b = \phi(a)$. Além disso, podemos incluir $Im(\phi)$ em $B$ fazendo $i : Im(\phi) \to B$ com $i(b) = b = \phi(a) = \theta(a + \ker \phi)$. Assim, \begin{equation}
        \phi(a) = i(\theta (a + \ker \phi)) = i(\theta(\pi(a)))
    \end{equation}
\end{demonstracao}

\end{document}